\chapter{Evaluation}
\label{cha:eval}

This chapter describes how I evaluated the laboratory and which numbers I can actually cite. I only use values from files that exist in the repository. If a file is missing, I write that the number is \textit{not yet exported}. I do not invent clean accuracy, robust accuracy, or attack success rate.

\section{What I evaluated}
\label{sec:eval-what}

I separated three kinds of output.

\begin{enumerate}
\item \textbf{Configuration.} Scores I assigned in templates. The dashboard 75\% and 51\% values belong here.
\item \textbf{Computed on a reconstructed corpus.} The phishing JSON in \texttt{eval/phishing/results/phishing\_metrics.json}.
\item \textbf{Computed on a real classifier.} The adversarial JSON in \texttt{eval/adversarial/results/adversarial\_metrics.json}. The file currently holds an 8-image debug probe. I do not cite those eight images as the 200-image run my supervisor asked for.
\end{enumerate}

I also ran the 25 unit tests and I read the \texttt{source} field in \texttt{server.ts}.

\section{Dashboard values are configuration}
\label{sec:eval-a3}

A3 in the supervisor action list is explicit. The 75\% and 51\% figures are not findings.

When I click ``Load Study Dataset (90 Samples)'', \texttt{generate90Samples()} fills realism and detectability from the bands in Section~\ref{sec:impl-phish}, using seed 20260927. The dashboard uses \texttt{Math.round} on those means. For that seeded load the rounded values are 75\% mean designed realism (raw mean $75.12$) and 51\% for $100$ minus mean AI-class detectability (raw $51.23$). Those are still configuration. They are not filter measurements.

I wrote those bands so that conventional rows look easy to spot and AI-style rows look harder. That is a teaching design. It is not a measurement of a mail filter, a clinical pilot, or a sample-size study. I deleted that kind of language from the user interface. I do not use it here.

The dataset is in-memory templates. It is not a mailbox harvest.

\section{Phishing benchmark}
\label{sec:eval-phish}

I wrote \texttt{eval/phishing/run\_benchmark.py}. The first corpus was 288 repeated templates and the classifier reached F1 $= 1.0$. That score was a construction error, not a result, so I rebuilt the set. Ham and phish now share the same stems and differ mainly in the last instruction. The sender domain is stored but held out of the TF-IDF text, so the model is a wording check like the teaching tab.

The current file has 187 reconstructed rows, split 70/30 with \texttt{random\_state=20260927}. Table~\ref{tab:phish} copies the exported numbers. I do not treat them as Gmail or SpamAssassin performance.

\begin{table}[htbp]
\centering
\caption{Values copied from \texttt{eval/phishing/results/phishing\_metrics.json} (generated at 2026-10-20T18:10:55Z).}
\label{tab:phish}
\begin{tabular}{lr}
\toprule
Field & Value \\
\midrule
\texttt{n\_total} & 187 \\
\texttt{n\_train} & 130 \\
\texttt{n\_test} & 57 \\
Model & TF-IDF (1--2 gram) + logistic regression \\
\texttt{source} & computed \\
Accuracy & 0.649 \\
Precision & 0.633 \\
Recall & 0.939 \\
F1 & 0.756 \\
AUC & 0.75 \\
Test support (ham / phish) & 24 / 33 \\
\bottomrule
\end{tabular}
\end{table}

The JSON \texttt{label\_note} says the labels are the corpus labels shipped in \texttt{corpus.jsonl}, not a mailbox harvest and not Gemini-as-judge. Sender domain was held out. I agree with that note.

Recall is high and ham recall is low (0.25 on 24 ham test rows). The wording model still over-flags lookalike institutional text. That is a limitation of a small reconstructed set, not a detector finding about the internet.

I exported \texttt{MANUAL\_REVIEW.csv} and I ticked 20 rows. The teaching label matched the last instruction in each of those rows (\texttt{rinor\_agrees=yes}).

\section{Adversarial method}
\label{sec:eval-adv}

I implemented the method my supervisor asked for: FGSM and PGD~\cite{goodfellow2015explaining,madry2018towards}, on at least 200 images, with clean accuracy, robust accuracy, attack success rate, and perturbation size. Torchattacks is used when the import works~\cite{kim2020torchattacks}.

\texttt{eval/adversarial/run\_attacks.py} does the following.

\begin{enumerate}
\item Download the CIFAR-10 test set~\cite{krizhevsky2009cifar}.
\item Draw $n=200$ indices with \texttt{numpy} RandomState seed 20260927.
\item Load \texttt{cifar10\_resnet56} from \texttt{chenyaofo/pytorch-cifar-models}~\cite{chenyaofo2021cifar}.
\item Measure clean accuracy.
\item For $\varepsilon \in \{2/255,\ 8/255,\ 16/255\}$, measure FGSM and PGD robust accuracy, attack success rate (label flip versus the clean prediction), and mean $L_\infty$.
\item Write \texttt{eval/adversarial/results/adversarial\_metrics.json} with \texttt{source: "computed"}.
\end{enumerate}

The JSON file exists, but it is an 8-image debug probe (\texttt{n=8}, clean accuracy $0.875$, seed 20260927). I used that run only to check that CIFAR normalisation was wired. I do not copy those eight-image FGSM/PGD scores into a results table. The supervisor asked for at least 200 images. Until \texttt{n=200} is in that file:

\begin{itemize}
\item clean accuracy on 200 images: \textit{not yet exported};
\item robust accuracy (FGSM, PGD, each $\varepsilon$) on 200 images: \textit{not yet exported};
\item attack success rate on 200 images: \textit{not yet exported};
\item mean $L_\infty$ on 200 images: \textit{not yet exported}.
\end{itemize}

The live Simulate button is a single-image call to \texttt{POST /attack}. It returns one clean class, one adversarial class, and two softmax percentages. Those values are model outputs when the service is up. They are not the batch summary. I do not copy a live button result into this chapter because I do not have a frozen export of that click.

\section{Unit tests}
\label{sec:eval-tests}

I ran \texttt{npm test}. Twenty-five tests passed in \texttt{tests/labSecurity.test.ts} and \texttt{tests/apiContract.test.ts}. They check string comparison, fail-closed credentials, the injection regex, the 10/40 rate limits, the 32\,kb body cap, and the four documented \texttt{source} labels. They do not measure phishing F1 or ASR.

\section{What I did not evaluate}
\label{sec:eval-not}

I did not run a user study. I did not measure Gmail catch rates. I did not score real videos. I did not compare ResNet-50 ImageNet with ResNet-56 CIFAR in a finished metrics file. I did not count live \texttt{source=gemini} versus \texttt{source=fallback} calls into a JSON export. Those items stay as future work or as viva demonstrations, not as hidden numbers.
